First take \(S=H_n\).  The score constraint is then automatic at every \(\rho\), so \(q_\rho(A;H_n)=\lambda_{\max}(A\mid H_n)\).  Every \(A\in\mathcal C_n\) has trace \(n\) on the \((n-1)\)-dimensional space \(H_n\), whence \(\lambda_{\max}(A\mid H_n)\geq c_n\).  Complete randomization has covariance \(c_nP_H\) and attains equality.  Thus \(v^*(\rho;H_n)=c_n\) for all \(\rho\), and no strict improvement is possible.

For \(n=4\), supplement \(u_4\) by
\[
v_4=\frac12(1,-1,1,-1)^\top,
\qquad w_4=\frac12(1,-1,-1,1)^\top.
\]
These three vectors form an orthonormal basis of \(H_4\).  Up to sign, the three balanced assignments are \(2u_4,2v_4,2w_4\).  Hence every \(A\in\mathcal C_4\) is diagonal in this basis with eigenvalues
\[
4p_u,\quad4p_v,\quad4p_w,
\qquad p_u,p_v,p_w\geq0,\quad p_u+p_v+p_w=1.
\]
Put \(s=\rho^2\) and \(r=\max(p_v,p_w)\).  Allocating the component orthogonal to \(u_4\) to the larger of the last two eigen-directions gives
\[
q_\rho(A;S_4)
=4\max_{a^2\in[s,1]}\{p_u a^2+r(1-a^2)\}
=4\max\{p_u,\ p_us+r(1-s)\}.
\]
For fixed \(p_u=p\), the smallest possible \(r\) is \((1-p)/2\), attained by \(p_v=p_w=(1-p)/2\).  Therefore
\[
v^*(\rho;S_4)
=4\min_{0\leq p\leq1}
\max\left\{p,\ ps+\frac{(1-p)(1-s)}2\right\}.
\]
If \(s\leq1/3\), the minimum occurs at \(p=1/3\) and equals \(4/3\).  Indeed the two affine terms meet there, and moving to either side cannot lower their maximum.  If \(s\geq1/3\), the second affine term is nondecreasing in \(p\), so \(p=0\) is optimal and the value is \(2(1-s)\).  This proves
\[
v^*(\rho;S_4)=\min\left\{\frac43,\,2(1-\rho^2)\right\}.
\]
The choice \(p_u=p_v=p_w=1/3\) is complete randomization.  The choice \(p_u=0\), \(p_v=p_w=1/2\) is precisely the displayed law \(d_{\mathrm{pair},4}\), and direct averaging gives its displayed covariance \(A_{\mathrm{pair},4}\).  This covariance annihilates \(u_4\), so \(\kappa(S_4)=0\).  The strict-improvement set from the formula is \(\rho^2>1/3\).  The admissibility-frontier theorem consequently gives
\[
\gamma(S_4)=\frac13,
\qquad \rho^*(S_4)=\frac1{\sqrt3}.
\]

For general even \(n\), assign one treated and one control unit within every fixed pair, with independent fair orientations.  Every supported assignment is balanced and the distribution has mean zero.  If \(x\in S_{\mathrm{pair}}\), then
\[
z^\top x=\sum_{j=1}^{n/2}x_{2j-1}(z_{2j-1}+z_{2j})=0.
\]
Thus its covariance annihilates \(S_{\mathrm{pair}}\), proving \(\kappa(S_{\mathrm{pair}})=0<c_n\) and hence \(S_{\mathrm{pair}}\in\mathfrak A_n\).  At \(n=8\), perfect quality restricts \(\mu\) to this annihilated space, so \(v^*(1;S_{\mathrm{pair}})=0\) and, for \(M=1\),
\[
V^*(1;S_{\mathrm{pair}})=\frac{4}{8}\,0=0,
\qquad V_{\mathrm{CRD}}(1;S_{\mathrm{pair}})=\frac4{8-1}=\frac47.
\]
The exhaustive three-sign-pair calculation above is also the requested smallest exact stress test: it checks every balanced assignment at \(n=4\), including both sides of the switching boundary.